When the endpoint tensors are injective, a matrix inserted on the bond $e$ is
realized by a physical operation on a neighbouring physical leg, and conversely.
This sets up an isomorphism between the bond-$e$ matrix algebra of $A$ and that
of $B$, which by the Skolem--Noether theorem is conjugation by an invertible
matrix $X_e$, the edge gauge: for every physical configuration $\sigma$ and bond
matrix $M$,
$C_A(e,\sigma,M)=C_B(e,\sigma,X_eMX_e^{-1})$.
Collecting the $X_e$ over all edges and reconciling them on the shared bonds
produces the global gauge family that relates the two tensors.

\begin{definition}[Finite kernel descent datum]%
    \label{def:peps_finiteRegionKernelDescent}
    \lean{TNLean.PEPS.FiniteRegionKernelDescent}
    \leanok
    Let $V$ be a vertex type with decidable equality. A finite kernel descent
    datum is a predicate $K$ on finite subsets $S\subseteq V$ together with
    implications
    $j\in S\Longrightarrow K(S)\Longrightarrow K(S\setminus\{j\})$.
    In the edge-blocking proof, $K(S)$ is the assertion that the boundary
    coefficients $c_{\rho}$, with the virtual indices exposed at that stage,
    give zero after the tensors in $S$ have been contracted.
\end{definition}

\begin{lemma}[Finite descent of kernel conditions]%
    \label{lem:peps_finiteRegionKernelDescent}
    \lean{TNLean.PEPS.FiniteRegionKernelDescent.descend_to_empty}
    \lean{TNLean.PEPS.FiniteRegionKernelDescent.descend}
    \leanok
    \uses{def:peps_finiteRegionKernelDescent}
    For a finite kernel descent datum and any finite $R\subseteq V$,
    $K(R)\Longrightarrow K(\varnothing)$. More generally, if
    $K(\varnothing)\Longrightarrow Q$, then $K(R)\Longrightarrow Q$.
\end{lemma}
\begin{proof}\leanok
    Induct on the finite set $R$. If $R=\varnothing$, there is nothing to
    prove. Otherwise write $R=S\cup\{j\}$ with $j\notin S$. The deletion
    implication gives $K(R)\Longrightarrow K(S)$, and the induction hypothesis
    gives $K(S)\Longrightarrow K(\varnothing)$.
\end{proof}

\begin{definition}[Edge-middle kernel descent data]%
    \label{def:peps_edgeMiddleKernelDescentData}
    \lean{TNLean.PEPS.EdgeMiddleBoundaryLabel,
        TNLean.PEPS.EdgeMiddleKernelDescentData}
    \leanok
    \uses{def:peps_edgeMiddleTensorFamily,
        def:peps_finiteRegionKernelDescent}
    Let $e=(u,v)$ and $M_e=V\setminus\{u,v\}$. An edge-middle kernel descent
    datum is indexed by boundary labels $\rho$ consisting of the two residual
    endpoint configurations. It assigns to every finitely supported coefficient
    family $c=(c_{\rho})_{\rho}$ a finite kernel descent datum $K_c(S)$ such that
    \begin{align}
        \left(\forall\tau,\
            \sum_{\rho}c_{\rho}T_{A,M_e}^{\rho}(\tau)=0\right)
        &\Longrightarrow K_c(M_e),
        \notag\\
        K_c(\varnothing)
        &\Longrightarrow \forall\rho,\ c_{\rho}=0.
        \notag
    \end{align}
\end{definition}

\begin{theorem}[Kernel descent gives edge-middle injectivity]%
    \label{thm:peps_edgeMiddleTensorInjective_of_kernelDescent}
    \lean{TNLean.PEPS.EdgeMiddleKernelDescentData.edgeMiddleTensorInjective,
        TNLean.PEPS.IsVertexInjective.edgeBlockedThreeSiteInjective_of_kernelDescent}
    \leanok
    \uses{def:peps_edgeMiddleKernelDescentData,
        lem:peps_finiteRegionKernelDescent,
        thm:peps_edgeBlockedThreeSiteInjective_of_middle}
    Let $e=(u,v)$ and $M_e=V\setminus\{u,v\}$. If an edge-middle kernel descent
    datum is given, then $T_{A,M_e}$ is injective. Consequently, if $A$ is
    vertex-injective, the edge-blocked three-site chain at $e$ is injective.
\end{theorem}
\begin{proof}\leanok
    Let $c=(c_{\rho})_{\rho}$ be a finite linear relation among the middle
    tensor vectors:
    $\sum_{\rho}c_{\rho}T_{A,M_e}^{\rho}(\tau)=0$ for every $\tau$.
    The initial part of the descent datum gives $K_c(M_e)$. The finite descent
    lemma gives $K_c(M_e)\Longrightarrow K_c(\varnothing)$, and the terminal
    part gives $\forall\rho,\ c_{\rho}=0$. Thus every finite relation is
    trivial, which is $\inj(T_{A,M_e})$. Vertex injectivity at
    $u$ and $v$ gives $\inj(T_{A,u})$ and
    $\inj(T_{A,v})$; the edge-blocked three-site injectivity at
    $e$ follows.
\end{proof}

\begin{theorem}[Edge blocking preserves injectivity]%
    \label{thm:peps_edgeBlockedThreeSiteInjective}
    \lean{TNLean.PEPS.IsVertexInjective.edgeBlockedThreeSiteInjective}
    \leanok
    \uses{def:IsVertexInjective, def:peps_edgeBlockedThreeSiteInjective,
        thm:peps_edgeMiddleTensorInjective_of_kernelDescent,
        lem:peps_localCoeff_eq_zero_of_contract_zero,
        lem:peps_not_edgeMiddleTensorInjective_of_isEmpty}
    For a tensor map or blocked tensor family $T$, write
    $\inj(T)$ for the assertion that $T$ is injective. If $A$
    is vertex-injective and every bond dimension is positive, then for every
    edge $e=(u,v)$,
    $\inj(T_{A,u})$, $\inj(T_{A,V\setminus\{u,v\}})$, and
    $\inj(T_{A,v})$.
    Here $T_{A,u}$ and $T_{A,v}$ are the two endpoint tensor maps, and
    $T_{A,V\setminus\{u,v\}}$ is the blocked middle tensor family. This is the
    precise assertion after the edge-blocking diagram in
    \cite[Section~3]{Molnar2018NormalPEPS}. The positive-bond hypothesis is the
    source's standing assumption that injective PEPS have nonzero virtual bond
    spaces. Without it, an empty complement configuration space can make the
    middle tensor vanish identically;
    Lemmas~\ref{lem:peps_edgeMiddleTensorFamily_eq_zero_of_isEmpty}
    and~\ref{lem:peps_not_edgeMiddleTensorInjective_of_isEmpty} record
    this obstruction.
\end{theorem}
\begin{proof}\leanok
    \uses{thm:peps_edgeMiddleTensorInjective_of_kernelDescent,
        lem:peps_localCoeff_eq_zero_of_contract_zero}
    Fix $e=(u,v)$ and put $M_e=V\setminus\{u,v\}$. From vertex injectivity,
    construct the edge-middle descent datum required by
    Theorem~\ref{thm:peps_edgeMiddleTensorInjective_of_kernelDescent}: for each
    finite relation
    $\sum_{\rho}c_{\rho}T_{A,M_e}^{\rho}(\tau)=0$ for every $\tau$,
    the associated conditions $K_c(S)$ satisfy $K_c(M_e)$,
    $j\in S\Longrightarrow K_c(S)\Longrightarrow K_c(S\setminus\{j\})$, and
    $K_c(\varnothing)\Longrightarrow \forall\rho,\ c_{\rho}=0$. The descent
    implication is the local contraction identity: for $j\in S$, vertex
    injectivity at $j$ gives a left inverse $L_j$, and applying $L_j$ to the
    relation contracts away that tensor,
    $j\in S,\ K_c(S)\mathrel{\stackrel{L_j}{\Longrightarrow}}
    K_c(S\setminus\{j\})$.
    The contraction at $j$ uses the factorization
    \begin{align}
        \mc C_e
        &\cong
        \mc V_j
        \times
        \prod_{\substack{f\ne e\\ f\not\ni j}}
        \{0,\ldots,\dim A_f-1\},
        \notag
    \end{align}
    where $\mc C_e$ is the complement virtual-configuration space for
    the distinguished edge $e$ and $\mc V_j$ is the local virtual
    configuration space at $j$. The terminal step
    $K_c(\varnothing)\Longrightarrow \forall\rho,\ c_{\rho}=0$ isolates each
    boundary coefficient by surjectivity of the boundary labelling, which is
    where the positive bond dimensions are used to fill the interior virtual
    indices of a witness configuration.
    The kernel descent theorem gives $\inj(T_{A,M_e})$. The same
    theorem, together with vertex injectivity at $u$ and $v$, gives
    $\inj(T_{A,u})$, $\inj(T_{A,M_e})$, and
    $\inj(T_{A,v})$ as the edge-blocked three-site injectivity
    assertion.
\end{proof}

\begin{theorem}[Local virtual insertion realized physically on either endpoint]%
    \label{thm:peps_virtualInsertionPhysicalRealization}
    \lean{TNLean.PEPS.edgeVirtualInsertionPhysicalRealization}
    \leanok
    \uses{def:IsVertexInjective, def:peps_localIncidentMatrixOp,
        thm:peps_localIncidentMatrixOp_physicalRealization}
    Let $A$ be vertex-injective and let $e=(u,v)$ be an edge. For every
    matrix $M$ on the bond space of $e$, the one-leg virtual action induced
    by $M$ on the local tensor at $u$, and likewise at $v$, is realized by
    a physical operator on that endpoint's physical space.
    On a closed three-site chain, the corresponding state equality is
    % Formula (paper_normal.tex:303--352): inserting M on the (1,2) bond equals
    % applying O_1 to the physical leg of A_1 and also equals applying O_2 to the
    % physical leg of A_2.  Ink-to-index: the M panel has four virtual joins,
    % while each O panel has the three virtual joins of the closed chain; their
    % three north legs are the physical indices, with O_1 or O_2 inserted on the
    % indicated physical leg.
    % Boundary signature: every panel has (V,P)=(0,3).  This is the exact
    % closed-chain consequence drawn in Papers/1804.04964/paper_normal.tex:303--330,
    % not the boundary signature of the theorem's general local tensor.
    \begin{center}
        \tenkzkernel
        \begin{tenkz}[rows={ket}, cols=3, physical=up, boundary=periodic]
            \tn{A_1} & \tn{A_2} & \tn{A_3}
            \tn[skin=ring, species=operator, at=on bond-1-1-1-2 0.5]{M}
        \end{tenkz}
        $ = $
        \begin{tenkz}[rows={ket}, cols=3, physical=up, boundary=periodic]
            \tn[name=opSite]{A_1} & \tn{A_2} & \tn{A_3}
            \tn[at=1 n of opSite, skin=ring, species=operator,
                name=opRing, ports={270:physical, 90:physical}]{O_1}
            \tnwire{opSite.90}{opRing.270}
        \end{tenkz}
        $ = $
        \begin{tenkz}[rows={ket}, cols=3, physical=up, boundary=periodic]
            \tn{A_1} & \tn[name=opSite]{A_2} & \tn{A_3}
            \tn[at=1 n of opSite, skin=ring, species=operator,
                name=opRing, ports={270:physical, 90:physical}]{O_2}
            \tnwire{opSite.90}{opRing.270}
        \end{tenkz}
    \end{center}

    % This is the local endpoint part of Molnar2018NormalPEPS, Section 3. The corresponding
    % coefficient identity is
    % thm:peps_edgeInsertedCoeff_endpointPhysicalRealization.
\end{theorem}
\begin{proof}\leanok
    Apply the one-edge realization theorem to the two endpoint incidences of
    the distinguished edge.
\end{proof}

\begin{theorem}[Left-endpoint projected recovery]%
    \label{thm:peps_edgeLeftLocalVirtualOpOfPhysicalOp_eq_iff_projected_realization_eq}
    \lean{TNLean.PEPS.edgeLeftLocalVirtualOpOfPhysicalOp_eq_iff_projected_realization_eq}
    \leanok
    \uses{thm:peps_localVirtualOpOfPhysicalOp_eq_iff_projected_realization_eq,
        def:peps_localIncidentMatrixOp}
    Let $e=(u,v)$. A physical operator at the left endpoint pulls back to
    the one-edge virtual action of $M^{\mathsf T}$ if and only if its
    compressed physical action is the canonical realization of that
    one-edge virtual action.
\end{theorem}
\begin{proof}\leanok
    This is the projected local recovery equivalence applied to the left
    incident edge of $e$, with the transpose convention used by the
    left-endpoint coefficient identity.
\end{proof}

\begin{theorem}[Right-endpoint projected recovery]%
    \label{thm:peps_edgeRightLocalVirtualOpOfPhysicalOp_eq_iff_projected_realization_eq}
    \lean{TNLean.PEPS.edgeRightLocalVirtualOpOfPhysicalOp_eq_iff_projected_realization_eq}
    \leanok
    \uses{thm:peps_localVirtualOpOfPhysicalOp_eq_iff_projected_realization_eq,
        def:peps_localIncidentMatrixOp}
    Let $e=(u,v)$. A physical operator at the right endpoint pulls back to
    the one-edge virtual action of $M$ if and only if its compressed physical
    action is the canonical realization of that one-edge virtual action.
\end{theorem}
\begin{proof}\leanok
    This is the projected local recovery equivalence applied to the right
    incident edge of $e$.
\end{proof}

\begin{theorem}[Endpoint recovery of a common virtual insertion]%
    \label{thm:peps_edgeEndpointLocalVirtualOpOfPhysicalOp_eq_of_projected_realization_eq}
    \lean{TNLean.PEPS.edgeEndpointLocalVirtualOpOfPhysicalOp_eq_of_projected_realization_eq}
    \leanok
    \uses{thm:peps_edgeLeftLocalVirtualOpOfPhysicalOp_eq_iff_projected_realization_eq,
        thm:peps_edgeRightLocalVirtualOpOfPhysicalOp_eq_iff_projected_realization_eq}
    Let $e=(u,v)$. If
    $P_uO_1P_u=\Phi_uT_{u,e}(M^{\mathsf T})\Psi_u$ and
    $P_vO_2P_v=\Phi_vT_{v,e}(M)\Psi_v$, then
    $\Psi_uO_1\Phi_u=T_{u,e}(M^{\mathsf T})$ and
    $\Psi_vO_2\Phi_v=T_{v,e}(M)$.
    Here $T_{w,e}(N)$ denotes the virtual operator on the legs incident to $w$
    which acts by $N$ on the leg belonging to $e$ and by the identity on the
    other incident legs.
\end{theorem}
\begin{proof}\leanok
    The endpoint equivalences give
    \begin{align}
        P_uO_1P_u=\Phi_uT_{u,e}(M^{\mathsf T})\Psi_u
        &\iff
        \Psi_uO_1\Phi_u=T_{u,e}(M^{\mathsf T}),
        \notag\\
        P_vO_2P_v=\Phi_vT_{v,e}(M)\Psi_v
        &\iff
        \Psi_vO_2\Phi_v=T_{v,e}(M).
        \notag
    \end{align}
    The two hypotheses are the left sides of these equivalences.
\end{proof}

\begin{theorem}[Endpoint physical-to-virtual recovery under image preservation]%
    \label{thm:peps_edgePhysicalToVirtualInsertion_of_projected_realization_eq}
    \lean{TNLean.PEPS.edgePhysicalToVirtualInsertion_of_projected_realization_eq}
    \leanok
    \uses{thm:peps_edgeEndpointLocalVirtualOpOfPhysicalOp_eq_of_projected_realization_eq,
        thm:peps_localVirtualOpOfPhysicalOp_realizes_of_projector}
    Let $e=(u,v)$. Suppose that
    $P_uO_1P_u=\Phi_uT_{u,e}(M^{\mathsf T})\Psi_u$ and
    $P_vO_2P_v=\Phi_vT_{v,e}(M)\Psi_v$. Suppose also that
    $O_1$ maps the image of $\Phi_u$ into itself and $O_2$ maps the image of
    $\Phi_v$ into itself. Then
    \begin{align}
        O_1\Phi_u
        &=\Phi_uT_{u,e}(M^{\mathsf T}),
        \notag\\
        O_2\Phi_v
        &=\Phi_vT_{v,e}(M).
        \notag
    \end{align}
\end{theorem}
\begin{proof}\leanok
    The preceding endpoint theorem identifies the two virtual pullbacks with
    the one-edge actions of $M^{\mathsf T}$ and $M$. The pullback
    specification gives the projected physical actions, and the
    image-preservation hypotheses remove the endpoint projectors.
\end{proof}

\begin{theorem}[Right-endpoint physical realization of the inserted coefficient]%
    \label{thm:peps_edgeInsertedCoeff_rightPhysicalRealization}
    \lean{TNLean.PEPS.edgeInsertedCoeff_eq_sum_right_physicalRealization}
    \leanok
    \uses{def:peps_edgeInsertedCoeff,
        def:peps_edgeBoundaryRightIndexEquivInsertedBoundaryConfig,
        thm:peps_localTensorMap_localIncidentMatrixOp_single}
    Let $e=(u,v)$ and let $\sigma$ be a physical configuration. If a physical
    operator on $v$ realizes the one-edge virtual matrix action on the
    distinguished incident edge, then the inserted-edge coefficient at
    $\sigma$ is obtained by leaving the left endpoint and open middle
    contraction explicit and applying that physical operator to the right
    endpoint tensor.
\end{theorem}
\begin{proof}\leanok
    Expand the physical realization on a basis local configuration at the
    right endpoint. The resulting sum over the right bond index is then
    reindexed together with the ordinary edge-boundary data.
\end{proof}

\begin{theorem}[Left-endpoint physical realization of the inserted coefficient]%
    \label{thm:peps_edgeInsertedCoeff_leftPhysicalRealization}
    \lean{TNLean.PEPS.edgeInsertedCoeff_eq_sum_left_physicalRealization}
    \leanok
    \uses{def:peps_edgeInsertedCoeff,
        def:peps_edgeBoundaryLeftIndexEquivInsertedBoundaryConfig,
        thm:peps_localTensorMap_localIncidentMatrixOp_single}
    Let $e=(u,v)$ and let $\sigma$ be a physical configuration. If a physical
    operator on $u$ realizes the one-edge virtual action of $M^{\mathsf T}$ on
    the distinguished incident edge, then the inserted-edge coefficient for
    $M$ at $\sigma$ is obtained by applying that physical operator to the left
    endpoint tensor and leaving the open middle contraction and the right
    endpoint explicit.
\end{theorem}
\begin{proof}\leanok
    Expand the physical realization on a basis local configuration at the left
    endpoint. The transpose converts the ordinary right bond index and the new
    left bond index into the coefficient $M_{ij}$; the resulting finite sum
    is then reindexed with the ordinary edge-boundary data.
\end{proof}

\begin{theorem}[Endpoint physical realizations of the inserted coefficient]%
    \label{thm:peps_edgeInsertedCoeff_endpointPhysicalRealization}
    \lean{TNLean.PEPS.edgeInsertedCoeff_endpointPhysicalRealization}
    \leanok
    \uses{thm:peps_localIncidentMatrixOp_physicalRealization,
        thm:peps_edgeInsertedCoeff_leftPhysicalRealization,
        thm:peps_edgeInsertedCoeff_rightPhysicalRealization}
    If the PEPS tensor is vertex-injective, then for every physical
    configuration $\sigma$, each inserted edge matrix is realized by physical
    operators at the endpoint tensors: the left endpoint realizes
    $M^{\mathsf T}$, the right endpoint realizes $M$, and both realizations
    reproduce the same inserted-edge coefficient at $\sigma$ in the
    edge-centered three-site decomposition.
\end{theorem}
\begin{proof}\leanok
    Apply the local physical-realization theorem at the two endpoints and then
    use the two coefficient identities above.
\end{proof}

\begin{theorem}[Equal neighboring physical actions recover a virtual insertion]%
    \label{thm:peps_physical_to_virtual_insertion}
    \lean{TNLean.PEPS.physical_to_virtual_insertion}
    \leanok
    \uses{def:peps_edgeBlockedThreeSiteInjective, def:peps_edgeBoundaryConfig,
        def:peps_edgeLeftLocalConfig, def:peps_edgeRightLocalConfig,
        def:peps_edgeOpenMiddleWeight, def:peps_localTensorMap,
        def:peps_localIncidentMatrixOp}
    Let $e=(u,v)$. Assume that the edge-blocked three-site chain at $e$ is
    injective and that every virtual bond has positive dimension. Let $O_1,O_2$
    be physical operators at the endpoint tensors, and
    write $\Phi_u,\Phi_v$ for the two local tensor maps. Suppose that, for every
    physical configuration $\sigma$,
    \begin{align}
        &\sum_{\beta}
        O_1(A_u(\beta_u))_{\sigma_u}
        C_{\mathrm{mid}}(\sigma;\beta_{\mathrm{L}},\beta_{\mathrm{R}})
        A_v(\beta_v)_{\sigma_v}
        \notag\\
        &\qquad =
        \sum_{\beta}
        A_u(\beta_u)_{\sigma_u}
        C_{\mathrm{mid}}(\sigma;\beta_{\mathrm{L}},\beta_{\mathrm{R}})
        O_2(A_v(\beta_v))_{\sigma_v}.
        \notag
    \end{align}
    Then
    \begin{align}
        \exists M\in\MN{D(e)},\qquad
        O_1\Phi_u
        &=\Phi_uT_{u,e}(M^{\mathsf T}),
        \notag\\
        O_2\Phi_v
        &=\Phi_vT_{v,e}(M).
        \notag
    \end{align}
    The same implication is represented diagrammatically by the
    $O_1,O_2\mapsto W$ step:
    % Formula (paper_normal.tex:355--485): O_1 acting at B_1 and O_2 acting at
    % B_2 give the same three-site state only if both arise from one bond matrix W.
    % Ink-to-index: each panel is the same closed virtual three-cycle; the first two
    % panels place O_1 and O_2 on physical legs, and the last splits the (1,2) bond
    % by W.  The three north legs remain the physical indices in every panel.
    % Boundary signature: each panel has (V,P)=(0,3).  Source verdict: exact
    % implication assembled from paper_normal.tex:355--375 and 422--485.
    \begin{center}
        \tenkzkernel
        \begin{tenkz}[rows={ket}, cols=3, physical=up, boundary=periodic]
            \tn[name=opSite]{B_1} & \tn{B_2} & \tn{B_3}
            \tn[at=1 n of opSite, skin=ring, species=operator,
                name=opRing, ports={270:physical, 90:physical}]{O_1}
            \tnwire{opSite.90}{opRing.270}
        \end{tenkz}
        $ = $
        \begin{tenkz}[rows={ket}, cols=3, physical=up, boundary=periodic]
            \tn{B_1} & \tn[name=opSite]{B_2} & \tn{B_3}
            \tn[at=1 n of opSite, skin=ring, species=operator,
                name=opRing, ports={270:physical, 90:physical}]{O_2}
            \tnwire{opSite.90}{opRing.270}
        \end{tenkz}
        $ \Longrightarrow $
        \begin{tenkz}[rows={ket}, cols=3, physical=up, boundary=periodic]
            \tn{B_1} & \tn{B_2} & \tn{B_3}
            \tn[skin=ring, species=operator, at=on bond-1-1-1-2 0.5]{W}
        \end{tenkz}
    \end{center}
    This is the converse part of the insertion-algebra lemma in
    \cite[Section~3]{Molnar2018NormalPEPS}, where the recovered matrix is
    denoted $W$.
\end{theorem}
\begin{proof}\leanok
    Contracting the middle block out of the hypothesis leaves a bond-contracted
    identity between the two endpoint tensors, holding for every residual
    boundary configuration. Inverting the right endpoint of this identity writes
    the action of $O_1$ on a left tensor vector as a bond-indexed combination of
    left tensor vectors with the same residual; the combining coefficients are
    read at a fixed reference right residual and define $M$. Inverting the left
    endpoint writes the action of $O_2$ on a right tensor vector as a bond-indexed
    combination of right tensor vectors. The left inverse identifies the two
    coefficient families, which is the source step $V=W$, so the right-endpoint
    combination is also governed by $M$. The reference residuals exist because
    every virtual bond has positive dimension.
\end{proof}

\begin{theorem}[Injectivity of an inserted edge coefficient]%
    \label{thm:peps_edgeInsertedCoeff_injective}
    \lean{TNLean.PEPS.edgeInsertedCoeff_injective}
    \leanok
    \uses{def:peps_edgeBlockedThreeSiteInjective, def:peps_edgeInsertedCoeff,
        thm:peps_edgeInsertedCoeff_leftPhysicalRealization,
        thm:peps_edgeInsertedCoeff_rightPhysicalRealization,
        thm:peps_physical_to_virtual_insertion}
    Let the edge-blocked three-site chain of $B$ at $e$ be injective, and
    assume every virtual bond of $B$ has positive dimension. If two matrices
    $M,M'\in\MN{D_B(e)}$ give the same inserted coefficient at every physical
    configuration, $C_B(e,\sigma,M)=C_B(e,\sigma,M')$ for every $\sigma$, then
    $M=M'$.
\end{theorem}

\begin{proof}\leanok
    Realize $M$ on the right endpoint and $M'^{\mathsf T}$ on the left
    endpoint. The endpoint realizations and the equality
    $C_B(e,\sigma,M)=C_B(e,\sigma,M')$ give, for every physical configuration
    $\sigma$,
    \begin{align}
        &\sum_{\beta}
        O_1(B_u(\beta_u))_{\sigma_u}
        C_{\mathrm{mid}}^B(\sigma;\beta_{\mathrm{L}},\beta_{\mathrm{R}})
        (B_v(\beta_v))_{\sigma_v}
        \notag\\
        &\qquad =
        \sum_{\beta}
        (B_u(\beta_u))_{\sigma_u}
        C_{\mathrm{mid}}^B(\sigma;\beta_{\mathrm{L}},\beta_{\mathrm{R}})
        O_2(B_v(\beta_v))_{\sigma_v}.
        \notag
    \end{align}
    Theorem~\ref{thm:peps_physical_to_virtual_insertion} gives a single virtual
    matrix realizing both endpoint operators. Endpoint uniqueness then
    identifies this matrix with $M$ and with $M'$.
\end{proof}
